\documentclass{article}

\usepackage{booktabs}
\usepackage{tabularx}

\title{Development Plan\\\progname}

\author{\authname}

\date{}

\input{../Comments.text}
\input{../Common.text}

\begin{document}

\maketitle

\begin{table}[hp]
\caption{Revision History} \label{TblRevisionHistory}
\begin{tabularx}{\textwidth}{llX}
\toprule
\textbf{Date} & \textbf{Developer(s)} & \textbf{Change}\\
\midrule
September 25th 2026 & Ayush Patel & Added intro and Sections 1 - 3 \\
September 26th 2026 & All Team Members & Added Team Meeting and Communication Plan and Team Roles \\
September 26th 2026 & All Team Members & Added Expected Technology and AI Use and Team Charter \\
Date2 & Name(s) & Description of changes\\
... & ... & ...\\
\bottomrule
\end{tabularx}
\end{table}

\newpage{}

\wss{Put your introductory blurb here.  Often the blurb is a brief roadmap of
what is contained in the report.}

\wss{Additional information on the development plan can be found in the
\href{https://gitlab.cas.mcmaster.ca/courses/capstone/-/blob/main/Lectures/L02b_POCAndDevPlan/POCAndDevPlan.pdf?ref_type=heads}
{lecture slides}.}

This document contains details about the development plan for the project.  It includes information about the
confidentiality and licensing (sections 1 - 3), team communication and meeting plans (sections 4 - 6), workflow and project decomposition (sections 7 - 8), proof of concept demonstration plan (section 9), expected technology (section 10), coding standard (section 11),
and appendices for generative AI use disclosure, reflection, and team charter.

\section{Confidential Information?}

\wss{State whether your project has confidential information from industry, or
not.  If there is confidential information, point to the agreement you have in
place.}

\wss{For most teams this section will just state that there is no confidential
information to protect.}

There is no confidential information to protect in this project.

\section{IP to Protect}

\wss{State whether there is IP to protect.  If there is, point to the agreement.
All students who are working on a project that requires an IP agreement are also
required to sign the ``Intellectual Property Guide Acknowledgement.''}


The intellectual property produced by this project consists solely of written content,
namely the code and documentation, which is covered under our copyright license, discussed below.

\section{Copyright License}

\wss{What copyright license is your team adopting.  Point to the license in your
repo.}

The copyright license adopted by the team for this project is the MIT license.
This license provides others the right to use, copy, modify, merge, publish, distribute, sublicense, and/or sell copies of the software,
provided that the copyright notice and permission notice are included in all copies.
The license also disclaims warranty and limits the authors' liability for any damages arising from use of the software.
The \href{https://github.com/Ayushpatel2026/Keystone/blob/main/LICENSE}{copyright license} for the project can be found in the team repository.

\section{Team Meeting Plan}

\wss{This section is for meetings between the team members, and meetings between 
the team and their supervisor/advisor (as appropriate).}

\wss{How often will the team meet? where? when?}

\wss{If the meeting is a physical location (not virtual), out of an abundance of
caution for safety reasons you shouldn't put the location online}

\wss{To help schedule meetings, you can use the lecture/tutorial time when we do 
not have anything else scheduled.}

\wss{How often will you meet with your supervisor/advisor?  when?  where?}

\wss{Will meetings be virtual?  At least some meetings should likely be
in-person.}

\wss{How will the meetings be structured?  There should be a chair for all meetings.  
There should be an agenda for all meetings.}


\begin{itemize}
    \item \textbf{Weekly Team Meeting:} Once a week for one hour -- Wednesdays, 2:30pm to 3:30pm, in person.
    Additional virtual meetings will be scheduled as needed when approaching a deliverable deadline.

    \item \textbf{Meeting Chair:} Waqar Ul-Hassan - prepares and sends out the meeting agenda and leads the meetings

    \item \textbf{Notetaker:} Ahmed Rashrash - record the meeting minutes using the template below

    \item \textbf{Supervisor Meetings:} We will meet with Dr. Smith at least once a month, and more frequently as needed. These will primarily occur during lecture time, with virtual meetings as a backup when an in-lecture meeting is not possible.

    \item \textbf{Agenda Requirement:} All meetings (team, TA, supervisor) must have an agenda, prepared and shared with the team at least 24 hours in advance by the chair. The agenda template is as follows:
    \begin{itemize}
        \item Date, time, attendees
        \item Purpose of the meeting
        \item Review of previous deliverable/meeting: what went well, what needs improvement
        \item Team member updates: completed tasks, outstanding tasks, tasks needing assistance
        \item Open questions for supervisor / advisor, TA, or stakeholder (as applicable)
        \item Action items and task owners for the next meeting
    \end{itemize}

    \item \textbf{Meeting Minutes Template:} Minutes will be recorded for every team, TA, and supervisor meeting, using the following structure:
    \begin{itemize}
        \item Date, time, attendees, absentees (with reason if provided)
        \item Summary of discussion per agenda item
        \item Decisions made
        \item Action items, owners, and due dates
        \item Open questions carried forward
    \end{itemize}
    Minutes will be stored in the description of a "Team Meeting" GitHub Issue created for each meeting, using the provided Team Meeting issue template.

\end{itemize}


\section{Team Communication Plan}

\wss{Issues on GitHub should be part of your communication plan.}

\begin{itemize}
    \item \textbf{Primary asynchronous channel:} Microsoft Teams chat, for day-to-day communication and quick questions.
    \item \textbf{Task-level communication:} GitHub -- issues, pull requests, and code review comments are used for anything related to a specific task, so discussion stays attached to the relevant work item. Members will \@-mention each other in PRs/issues, assign issues directly, and request reviews through GitHub's review request feature.
    \item \textbf{Meeting minutes:} Recorded and stored on GitHub (see Team Meeting Plan).
    \item \textbf{Supervisor/stakeholder communication:} Email (Dr. Smith prefers email), routed through the team liaison to keep a single point of contact.
    \item \textbf{Response-time expectation:} Each member will respond to direct communication in the team's Teams chat within 24 hours.
\end{itemize}

\section{Team Member Roles}

\wss{You should identify the types of roles you anticipate, like notetaker,
leader, meeting chair, reviewer.  Assigning specific people to those roles is
not necessary at this stage.  In a student team the role of the individuals will
likely change throughout the year.}


\begin{itemize}
    \item \textbf{Project-Level Roles:}
    \begin{itemize}
        \item Team Liaison: Ayush Patel -- primary point of contact with the supervisor and course staff.
        \item Meeting Chair: Waqar Ul-Hassan -- prepares and sends out the meeting agenda and leads the meetings.
        \item Notetaker: Ahmed Rashrash -- record the meeting minutes using the template below.
        \item GitHub Manager: Andrew Habib -- responsible for keeping labels, issue templates, branch protection rules, and the project board consistent and up to date.
    \end{itemize}

    \item \textbf{Documentation Quality Assurance:} Rotates per deliverable alphabetically by first name starting with Ahmed for the first deliverable.
    The member(s) in this role are responsible for checking the completed deliverable against both the repository checklist and the Avenue rubric before submission, and flagging any gaps to the team with enough time to address them before the deadline.
    \item This is in addition to, not a replacement for, normal PR review.

    \item \textbf{Technical Roles:}
    \begin{itemize}
        \item Frontend developer: Responsible for development of a user friendly interface - All members
        \item Backend developer: Responsible for developing the backend and databse layer - All members
        \item Reviewer (code): all members.
        \item Reviewer (documentation): all members.
        \item Technical QA: Responsible for doing functional testing after each code functionality related PR is merged. This will rotate for each deliverable and will be the same team member who is also the documentation QA.
    \end{itemize}
    We are intentionally keeping technical roles broad and overlapping at this stage rather than assigning rigid specializations, since the team is still early in the project and all members should be able to work across the stack.
    Specializations may emerge naturally as the project progresses.

    \item \textbf{Flexibility and Role Changes:} Roles are not fixed for the duration of the project. At the start of each milestone, the team will briefly revisit role assignments based on interest, workload balance, and skills members want to develop, and reassign as needed. This ensures no single member is locked into a role that no longer fits the project's needs or their own goals.

    \item \textbf{Overlap by Design:} Technical and review roles are deliberately assigned to "all" members rather than single individuals, so there is no single point of failure if one member is unavailable during a critical period.
\end{itemize}


\section{Workflow Plan}

\begin{itemize}
	\item How will you be using git, including branches, pull request, etc.?
	\item How will you be managing issues, including template issues, issue
	classification, etc.?
  \item Use of CI/CD
\end{itemize}

\section{Project Decomposition and Scheduling}

\begin{itemize}
  \item How will you be using GitHub projects?
  \item Include a link to your GitHub project
\end{itemize}

\wss{How will the project be scheduled?  This is the big picture schedule, not
details. You will need to reproduce information that is in the course outline
for deadlines.}

\section{Proof of Concept Demonstration Plan}

What is the main risk, or risks, for the success of your project?  What will you
demonstrate during your proof of concept demonstration to convince yourself that
you will be able to overcome this risk?

\section{Expected Technology}

\wss{What programming language or languages do you expect to use?  What external
libraries?  What frameworks?  What technologies.  Are there major components of
the implementation that you expect you will implement, despite the existence of
libraries that provide the required functionality.  For projects with machine
learning, will you use pre-trained models, or be training your own model?  }

\wss{The implementation decisions can, and likely will, change over the course
of the project.  The initial documentation should be written in an abstract way;
it should be agnostic of the implementation choices, unless the implementation
choices are project constraints.  However, recording our initial thoughts on
implementation helps understand the challenge level and feasibility of a
project.  It may also help with early identification of areas where project
members will need to augment their training.}

\begin{itemize}
    \item \textbf{Version Control and Project Management:} Git, Github and GitHub projects
    \item \textbf{Frontend:} React -- chosen for its large ecosystem, team familiarity, and strong component-based structure suited to an interactive web application. \textbf{Trade-off:} steeper initial setup compared to simpler frameworks, but this is offset by better long-term maintainability.
    \item \textbf{Backend:} Python with FastAPI -- chosen for fast development speed, built-in data validation, automatic API documentation, and strong support for any data-processing or ML-adjacent work the project may require. \textbf{Trade-off:} lower raw performance than a compiled-language backend, which is an acceptable trade-off given the project's expected scale.
    \item \textbf{Backend unit testing:} pytest -- chosen over the built-in \texttt{unittest} module for its more concise syntax, fixture system, and better plugin ecosystem for coverage reporting.
    \item \textbf{Frontend unit testing:} Jest -- chosen for its tight integration with React and minimal configuration overhead.
    \item \textbf{E2E Testing:} Playwright -- industry standard for end to end testing of web applications
    \item \textbf{Frontend linter:} ESLint -- industry standard for TypeScript/JavaScript, configurable to enforce team-agreed style rules.
    \item \textbf{Backend linter:} Ruff -- industry standard linter, easiest to setup, fastest and most modern
    \item \textbf{Unit test coverage tool:}  SonarQube -- industry standard for reporting and visualizing code coverage
    \item \textbf{Database:} PostgreSQL - relational database will make it easier
    \item \textbf{Application Packaging:} Docker - production standard to package application components and dependencies into one executable container \textbf{Trade-off:} Extra set up is required but this will ensure our code will run anywhere we choose to deploy in the future
    \item \textbf{Continuous Integration:} GitHub Actions. CI will run on every push and pull request, checking: linting (ESLint/Pylint), unit test pass/fail (Jest/pytest), and SonarQube coverage reports.
    \item \textit{Note: the technologies listed here reflect the team's current expectations and are not binding constraints. They may evolve as the project's design becomes more concrete.}
\end{itemize}

\section{Use of AI and LLMs}

\wss{This section is a continuation of the previous section.  The role of LLMs
in your project as important enough to warrant their own section.  
How will your team be using LLMs for code and documentation development?  
What tools will you be using?  How will you verify the output of your LLMs?}


\begin{itemize}
    \item \textbf{Purpose:} AI will be used to accelerate development and ensure deliverables align with the rubric and checklists, while the team retains ownership of all thinking, design, and logic decisions.
    \item \textbf{Tools:} Claude and ChatGPT for documentation and code implementation assistance, GitHub Copilot or Claude Code for in-editor code assistance.
    \item \textbf{Documentation use of AI:}
    \begin{itemize}
        \item Brainstorming and idea generation
        \item Proofreading, grammar, and spelling checks.
        \item Assistance with LaTeX syntax and compilation issues.
        \item Checking documents against the rubric and checklist for completeness.
    \end{itemize}
    Documentation content itself will be written by the team, not generated wholesale by AI.
    \item \textbf{Code use of AI:}
    \begin{itemize}
        \item A walking skeleton of the codebase must be written manually by the team.
        \item Function logic should first be written out as pseudocode by a team member; AI may then be used to help implement that pseudocode, rather than generating logic from scratch.
        \item AI is not used to make architecture or design decisions -- these remain human decisions.
        \item All AI-generated code must be reviewed line by line before being committed. Data structure choices, naming conventions, and other design decisions are made by the team, not the AI.
        \item AI may assist in generating unit tests and identifying edge cases, but all generated tests must be reviewed and understood by the team, not accepted blindly.
        \item Team members must manually verify that AI-generated tests are actually testing the intended behaviour (e.g. by intentionally breaking the implementation and confirming the test fails).
    \end{itemize}
    \item \textit{Note: this section covers AI's role in building the project. AI use in drafting this document itself is disclosed separately in the Generative AI Use Disclosure appendix.}
\end{itemize}


\section{Coding Standard}

\wss{What coding standard will you adopt?}

\newpage{}

\section{Appendix --- Generative AI Use Disclosure}

\wss{ChatGPT (version 5) was used to brainstorm ideas for the template for this disclose section.}

\wss{This section is about the use of AI to write this document, not your
planned use of AI for your project.  You discuss that topic in one of the
sections in the body of this report.}

\subsection{AI Tools Used}

\wss{Give the name of the AI tool(s) used and the version.}

\subsection{How and Where Used}

\wss{Explain what the tool(s) were used for.  Examples include: brainstorming,
explaining a technical concept, reviewing the document, generating or modifying
text, generating or modifying diagrams, identifying errors or omissions or
inconsistencies.}

\wss{For each of the uses, identify which parts of the document were affected.}

\wss{You do not need to report every interaction, but you should disclose the 
uses that materially contributed to your work}

\subsection{Verification of AI-Generated Content}

\wss{\begin{itemize}
    \item Describe how the team checked AI-generated or AI-assisted material
    before including it in the document.
    \item In particular, verify factual claims, technical assertions,
    requirements, calculations, references, and other information for which an
    AI system may produce incorrect or fabricated results.
    \item \textbf{The team is responsible for the accuracy, completeness,
consistency, and appropriateness of the submitted document, regardless of
whether material was generated by a human or an AI system.}
\end{itemize}}

\newpage{}

\section*{Appendix --- Reflection}

\wss{Not required for CAS 741}

\input{../Reflection.text}

\textbf{Waqar Ul-Hassan:}

\begin{enumerate}
    \item Why is it important to create a development plan prior to starting the
    project?
    
    For a project of this size, especially because we are developing it from scratch it is critical that we create a plan to ensure we have a scope of the project. We can set realistic expectations and everyone in the group is on the same page to ensure we don't have any miscommunication. Additionally, this way we are also able to identify risks earlier on by evaluating the project and thinking about it from a more management perspective.

    \item In your opinion, what are the advantages and disadvantages of using
    CI/CD?

    I feel that the biggest advantage of using CI/CD is to have an automated testing suite which will build when you push code and catch bugs, code formatting, etc. This ensure that we are keeping a certain standard of work in our active repository and overall are more confident when pushing to the main branch. The biggest disadvantage I think is definitely the maintenance of the CI/CD pipeline, making sure its up to date. On top of this if the pipeline is not made very will, it can cause some very frustrating problems and become a huge waste of time, especially when trying to scale a project. 

    \item What disagreements did your group have in this deliverable, if any,
    and how did you resolve them?

    Initially we had some disagreements on things like, which backend language we should with our react frontend. One of my group members wanted to use nextJS while others wanted to use python FastAPI. Another, disagreement we had was when deciding on the metric to use to check for contribution amount to the project. Some wanted to use lines of code, commits, and others wanted to use number of issues. For all these initial disagreements, we just held discussions during our meetings. In these discussions, we weighed the pros and cons of using each of the different methodologies. After weighing them together, we then decided for example to use python FastAPI as we feel we have more experience with it and also its much more flexible as well as use the number of issues as our comparison metric.
\end{enumerate}

\newpage{}

\section*{Appendix --- Team Charter}

\wss{borrows from
\href{https://engineering.up.edu/industry_partnerships/files/team-charter.pdf}
{University of Portland Team Charter}}

\subsection*{External Goals}

\wss{What are your team's external goals for this project? These are not the
goals related to the functionality or quality fo the project.  These are the
goals on what the team wishes to achieve with the project.  Potential goals are
to win a prize at the Capstone EXPO, or to have something to talk about in
interviews, or to get an A+, etc.}

\begin{itemize}
    \item Produce something we are proud of and can discuss meaningfully in interviews.
    \item Given the nature of the project, aim for it to be genuinely useful to future capstone students and course staff after we graduate.
    \item Improve our project management, design, and technical skills by applying them to a real-world problem.
    \item Get an A+
\end{itemize}

\subsection*{Attendance}

\subsubsection*{Expectations}

\wss{What are your team's expectations regarding meeting attendance (being on
time, leaving early, missing meetings, etc.)?}


\begin{itemize}
    \item 80\% attendance is required for team, TA, and supervisor meetings.
    \item If a member must miss a meeting, they must notify the team in advance and provide a written update covering their progress and any action items relevant to them.
    \item Members attending a meeting are expected to arrive on time (at most 10 minutes late) and stay for its full duration.
    \item \textbf{Trigger:} Missing 2 consecutive meetings, \emph{or} being late thrice in a row, triggers a consequence.
    \item \textbf{Consequence:} the member brings snacks/coffee to the next team meeting for a first occurrence; a second occurrence within the same term triggers a mandatory check-in meeting with the TA to discuss the pattern.
    \item These standards apply to non-emergency cases only (see below for emergencies).
\end{itemize}


\subsubsection*{Acceptable Excuse}

\wss{What constitutes an acceptable excuse for missing a meeting or a deadline?
What types of excuses will not be considered acceptable?}

\begin{itemize}
    \item \textbf{Acceptable:} midterm preparation, conflicting lectures/tutorials, work-related obligations, medical appointments.
    \item \textbf{Not acceptable:} forgetting the meeting, non-urgent personal plans, or repeated last-minute excuses of the same kind. Repeated use of the same excuse will be treated as a pattern and addressed under the "Stay on Track" consequences below.
    \item Advance notice is expected wherever the excuse is foreseeable (e.g. a scheduled appointment).
\end{itemize}


\subsubsection*{In Case of Emergency}

\wss{What process will team members follow if they have an emergency and cannot
attend a team meeting or complete their individual work promised for a team
deliverable?}

\begin{itemize}
    \item In a genuine emergency, the attendance rules above do not apply.
    \item The member should notify the team as soon as reasonably possible.
    \item If the emergency will materially impact a deliverable, there are two options:
    \begin{enumerate}
        \item The other members will redistribute the member's tasks among themselves, and then that team member can pick up more work in future deliverables.
        \item If the rest of the team cannot handle the extra workload, the team will raise it with the TA and/or course instructor proactively, rather than waiting until the deadline is missed.
    \end{enumerate}
\end{itemize}

\subsection*{Accountability and Teamwork}

\subsubsection*{Quality}

\wss{What are your team's expectations regarding the quality
of team members' preparation for team meetings and the quality of the
deliverables that members bring to the team?}

\begin{itemize}
    \item Deliverables must follow the AI usage guidelines in this document -- no unreviewed "AI slop."
    \item Every deliverable is checked against both the Avenue rubric and the repository checklist before submission, by the assigned documentation QA role for that deliverable (see Team Member Roles) in addition to the author.
    \item Members are expected to arrive at meetings ready to give a meaningful update or raise questions.
    \item Members should stay reasonably current on lecture content and upcoming deliverable requirements.
\end{itemize}


\subsubsection*{Attitude}

\wss{What are your team's expectations regarding team members' ideas,
interactions with the team, cooperation, attitudes, and anything else regarding
team member contributions?  Do you want to introduce a code of conduct?  Do you
want a conflict resolution plan?  Can adopt existing codes of conduct.}


\begin{itemize}
    \item Team members are expected to be positive, respectful, and willing to work hard.
    \item All members' ideas are given fair consideration; disagreement is handled respectfully, not dismissively.
    \item Members are expected to communicate openly rather than let frustration build silently.
\end{itemize}


\subsubsection*{Stay on Track}

\wss{What methods will be used to keep the team on track? How will your team
ensure that members contribute as expected to the team and that the team
performs as expected? How will your team reward members who do well and manage
members whose performance is below expectations?  What are the consequences for
someone not contributing their fair share?}

\wss{You may wish to use the project management metrics collected for the TA and
instructor for this.}

\wss{You can set target metrics for attendance, commits, etc.  What are the
consequences if someone doesn't hit their targets?  Do they need to bring the
coffee to the next team meeting?  Does the team need to make an appointment with
their TA, or the instructor?  Are there incentives for reaching targets early?}


\begin{itemize}
    \item Progress is tracked via the GitHub Project board -- issue status, assignee, and milestone fields are kept current so contribution is visible to the whole team, not just self-reported.
    \item \textbf{Regular Contributions:} Each member should have visible, regular activity (commits, PRs, or issue updates) across each week of an active deliverable period, not concentrated entirely in the final 1--2 days.
    \item \textbf{Target metrics:} Issues will be assigned at the beginning of each deliverable or during each weekly meeting. We will try to assign each person a relatively equal amount of work and
    finishing those issues within the provided deadline will be the main target metric.
    \item \textbf{Consequence for falling behind and repeated excuses:} First instance -- a direct, informal conversation with the member to understand blockers and rebalance tasks; a repeated pattern across multiple deliverables will result in a joint conversation involving the TA.
\end{itemize}


\subsubsection*{Team Building}

\wss{How will you build team cohesion (fun time, group rituals, etc.)? }

We will play ping pong together.

\subsubsection*{Decision Making}

\wss{How will you make decisions in your group? Consensus?  Vote? How will you
handle disagreements? }

\begin{itemize}
    \item Decisions are made by consensus where possible.
    \item Where consensus cannot be reached, the team will vote; the majority decision stands.
    \item Where a disagreement cannot be resolved by voting, the team will default to whichever option best aligns with the rubric and industry best practice; if still unresolved, the TA or supervisor will be brought in to mediate.
\end{itemize}

\end{document}